\section{The distributional current of the moving chart window}
\label{sec:v72-window-current}
The radial balance in Section~\ref{sec:v71-two-sided-current} is
an identity of functions at a common roof. To differentiate that
identity one must also differentiate the active window. We construct
the resulting distribution, including its entrance and exit terms.
The construction does not assume finite variation at the outer edge
of a selected disk.

Fix $m,R,\lambda=(n,k),\chi,\varepsilon$, and write $h=h_\chi$,
$j=h/2$, $L_z=L_{z,\lambda}$ and
$\overline L_z=j^{-1}\int_0^jL_z(v)\,dv$. All sums below run over
the same finite selected family at this fixed exact record. Translation
of a measure or distribution by $a$ is denoted by $\tau_a$.
For $j<H<h$ put
\begin{equation}\label{eq:v72-truncated-capacities}
 \begin{split}
 X_H(t)&=\sum_zJ_z\one_{(0,H)}(t-t_z)L_z(t-t_z),\\
 Y_H(t)&=\sum_zJ_z\overline L_z\one_{(0,H)}(t-t_z),
 \qquad C_H=X_H-Y_H.
 \end{split}
\end{equation}
The cutoff is an auxiliary exhaustion, not a deletion from the
physical source. The full functions $X,Y,C$ use $H=h$ and agree
with $\mathsf X,\mathsf Y,\mathsf C$ of the preceding section.

\begin{theorem}[Entrance and exit currents]
\label{thm:v72-window-current}
Let $H\in(j,h)$ be a continuity point of every $L_z$. The following
identities hold as finite signed Radon measures on $\R$:
\begin{align}
 D X_H={}&\sum_zJ_z\bigl\{
  \tau_{t_z}(DL_z|_{(0,H)})
       +L_z(0+)\delta_{t_z}-L_z(H-)\delta_{t_z+H}\bigr\},
                                  \label{eq:v72-DXH}\\
 D Y_H={}&\sum_zJ_z\overline L_z
                         (\delta_{t_z}-\delta_{t_z+H}),
                                  \label{eq:v72-DYH}\\
 D C_H={}&\sum_zJ_z\bigl\{
  \tau_{t_z}(DL_z|_{(0,H)})
       +(L_z(0+)-\overline L_z)\delta_{t_z}\notag\\
 &\hspace{39mm}-(L_z(H-)-\overline L_z)\delta_{t_z+H}\bigr\}.
                                  \label{eq:v72-DCH}
\end{align}
Each of the three measures has total mass zero. As $H\uparrow h$
through such values, the functions converge in $L^1(\R)$ to $X,Y,C$,
and their derivatives converge in distributions to $DX,DY,DC$.
In particular $DX$ and $DC$ are well-defined compactly supported
distributions of order at most one. They are not asserted to be
finite measures without an additional outer-edge variation bound.
\end{theorem}
\begin{proof}
Theorem~\ref{thm:v70-physical-current} gives $L_z\in BV(0,H')$
for every $H'<h$, including finite variation near zero. Thus the
traces in the assertion exist. There are only countably many
excluded values of $H$, even after the finite sum over $z$.
For $\phi\in C_c^\infty(\R)$, integration by parts for $BV$
functions gives
\[
 -\int_0^H L_z(u)\phi'(t_z+u)\,du
 =\int_{(0,H)}\phi(t_z+u)\,dDL_z(u)
  +L_z(0+)\phi(t_z)-L_z(H-)\phi(t_z+H).
\]
Multiply by $J_z$ and sum. A constant profile gives the second
formula, and subtraction gives the third. Testing against a smooth
function equal to one on the common compact support proves zero
total mass. Equivalently, the interior mass is
$L_z(H-)-L_z(0+)$ and cancels the two traces.

Since $0\le L_z,\overline L_z\le1$,
\[
 \|X-X_H\|_1+\|Y-Y_H\|_1
                  \le2(h-H)\sum_zJ_z.
\]
The sum is finite at the fixed count. The same bound holds for
$\|C-C_H\|_1$. Distributional differentiation is continuous under
$L^1$ convergence: the error against $\phi$ is bounded by the
corresponding $L^1$ error times $\|\phi'\|_\infty$. This also
proves independence of the exhausting sequence. No collision-uniform
bound on $\sum_zJ_z$ is used or obtained.
\end{proof}

The exit terms in \eqref{eq:v72-DCH} are indispensable. For example,
if $L_z=\one_{(a,h)}$ with $j<a<H$, then $\overline L_z=0$
and this chart contributes
$J_z(\delta_{t_z+a}-\delta_{t_z+H})$. Retaining only the interior
birth would give a measure of nonzero total mass, which cannot be
the derivative of this compactly supported function. This is a scalar
check of the boundary formula, not a claim about realizable Lorentz
words. If the outer trace does not exist, the complete right side of
\eqref{eq:v72-DCH} has a distributional limit; the exit terms need
not have a separate measure limit.

\begin{proposition}[Gluing and Fourier representation]
\label{prop:v72-window-fourier}
Artificial cuts of one profile cancel with their orientations before
any variation is taken. At a genuine jump the surviving atom is its
actual jump. With the Fourier convention
$\widehat f(\xi)=\int e^{-i\xi t}f(t)\,dt$, one has
\begin{equation}\label{eq:v72-window-fourier}
 \begin{split}
 \widehat{DX_H}(\xi)=\sum_zJ_ze^{-i\xi t_z}
 \bigg\{\int_{(0,H)}e^{-i\xi u}\,dDL_z(u)
       +L_z(0+)-e^{-i\xi H}L_z(H-)\bigg\}
       =i\xi\widehat X_H(\xi).
 \end{split}
\end{equation}
The analogous identities hold for $Y_H,C_H$ and for their full
window limits as distributions. For $\xi\ne0$ the right side can
be divided by $i\xi$; at zero the Fourier transform of the density,
not that quotient as an undefined expression, is used.
\end{proposition}
\begin{proof}
Split $(0,H)$ at $a$. The left piece has exit trace $-L_z(a-)$;
the right piece has entrance trace $L_z(a+)$. Their sum is the
actual jump. The interior derivatives of the two open pieces exclude
$a$, so this jump is counted once. At continuity it is zero. This
proves the gluing assertion. Test \eqref{eq:v72-DXH} against
$e^{-i\xi t}$ on the compact support, or apply the Fourier transform
to the distributional derivative. Passage to the full window is
justified by the $L^1$ convergence above. In particular
$|\widehat{DX}(\xi)|\le |\xi|\|X\|_1$ at each fixed count.
This is a finite-count bound, not a local-limit estimate.
\end{proof}

One must not identify $DC$ with the old scalar function $C$. Nor
may the present identities be called an operator-resolvent bound.
They supply the global derivative which such an argument would have
to use, with every active-window term retained. All translations and
physical coefficients occur before combining signed measures. No
cancellation between different $m,R,n,k$ is permitted.
